\documentclass[12pt]{amsart}

\usepackage{amsmath, amssymb, amsthm}
\usepackage{hyperref}
\usepackage{microtype}

\newtheorem{theorem}{Theorem}[section]
\newtheorem{lemma}[theorem]{Lemma}
\newtheorem{corollary}[theorem]{Corollary}
\newtheorem{proposition}[theorem]{Proposition}
\theoremstyle{definition}
\newtheorem{definition}[theorem]{Definition}
\newtheorem{remark}[theorem]{Remark}
\newtheorem{claim}[theorem]{Claim}
\newtheorem{conjecture}[theorem]{Conjecture}


\newcommand{\D}{\mathbb{D}}
\newcommand{\Z}{\mathbb{Z}}
\newcommand{\Zi}{\mathbb{Z}[i]}
\newcommand{\C}{\mathbb{C}}
\newcommand{\R}{\mathbb{R}}
\newcommand{\abs}[1]{\left|#1\right|}
\newcommand{\OD}{\mathcal{O}(\D)}
\newcommand{\mR}{\mathcal{R}}
\newcommand{\mRR}{\mathcal{R}_{\mathbb{R}}}

\title[Prime Spectra of Integer-Coefficient Power Series Rings]{The Prime Spectrum of Rings of\\Integer-Coefficient Power Series on the Disk}
\author{Jon Bannon}
\address{Department of Mathematics, Siena College, Loudonville, NY 12211, USA}
\author{David Victor Feldman}
\address{Department of Mathematics and Statistics, University of New Hampshire,
Durham, NH 03824, USA}
\date{August 2026}
\subjclass[2020]{Primary 30H50; Secondary 13F25, 13A15, 03C20, 68V20}
\keywords{Integer-coefficient power series, prime spectrum, maximal ideals,
ultraproducts, $p$-adic evaluation, unit disk, formal verification}

\begin{document}

\maketitle

\begin{abstract}
This is a sequel to \cite{PaperI}, where it is shown that a discrete
effective divisor on the open unit disk $\D$ is the zero divisor of a
holomorphic function with integer Taylor coefficients if and only if it
is invariant under complex conjugation.  Building on that realization
theorem and on the unit characterization established there, we study the
prime and maximal spectra of the rings $\mR = \Zi[[z]] \cap \OD$ and
$\mRR = \Z[[z]] \cap \OD$.  We give a trichotomy for maximal ideals and
classify the prime ideals containing $z$ or a nonzero integer: they are
$(z)$, the maximal ideals $(z,p)$, and the non-maximal primes $(p)$, for
$p$ a rational prime.  The remaining primes include, besides the class $\mathfrak{P}_1$
of primes containing a unit-constant-term element, a family of $p$-adic
evaluation primes $Q_p$ with $\mRR/Q_p \cong \Z_p$; in particular
$\mathrm{Spec}(\mRR) \setminus \mathfrak{P}_1$ strictly contains
$\{(0),(z)\} \cup \{(p) : p \text{ prime}\} \cup \{(z,p) : p
\text{ prime}\}$.  We
attach to each element of $\mathfrak{P}_1$ an analytic invariant
satisfying a partition property, and construct an injection from
ultrafilters on real-supported admissible divisors into $\mathfrak{P}_1$
by an ultraproduct device.  We prove the point-evaluation ideals $P_a$
are prime for all $a \in \D$ and maximal for real nonzero $a$, and deduce
a conditional B\'ezout property for functions with disjoint divisors and
coprime constant terms.  All results recalled from \cite{PaperI} are used
only as cited; no proof of that paper is reproduced here.
\end{abstract}

\section{Introduction}
\label{sec:intro}

The companion paper \cite{PaperI} proves that the only arithmetic
obstruction to prescribing the zero divisor of an integer-coefficient
holomorphic function on the unit disk $\D = \{z \in \C : \abs{z} < 1\}$
is conjugation invariance, and derives a first layer of ring-theoretic
consequences for the coefficient-restricted algebras
\[
  \mR = \Zi[[z]] \cap \OD, \qquad \mRR = \Z[[z]] \cap \OD,
\]
where $\Zi = \Z[i]$.  Those foundational results---the realization
theorems, the characterization of units, the factorization of
$\OD$-elements up to units, and the injectivity of the contraction
$\mathrm{MaxSpec}(\OD) \to \mathrm{Spec}(\mR)$---have been formally
verified.  We take them as given (Section~\ref{sec:recall}) and develop
from them a systematic description of the prime and maximal spectra.

The present paper contains no overlap with \cite{PaperI}: the
construction of realizing functions, its convergence analysis, and the
elementary-factor machinery live entirely in the companion paper and are
invoked here only through the statements collected in
Section~\ref{sec:recall}.  Everything from Section~\ref{sec:trichotomy}
onward is new.

\medskip\noindent\textbf{Results.}
After recalling the needed facts, we establish a trichotomy for maximal
ideals of $\mRR$ and show one of the three types cannot occur
(Section~\ref{sec:trichotomy}).  We then introduce the three ambient
classes of prime ideals $\mathfrak{P} \supseteq \mathfrak{P}_1 \supseteq
\mathfrak{M} \cap \mathfrak{P}_1$ and classify the primes containing $z$
or a nonzero integer: they are $(z)$ and the $(z,p)$ (the latter maximal)
when $z \in P$, and the $(p)$ (prime, non-maximal) when $z \notin P$ and
$P \cap \Z \neq 0$ (Section~\ref{sec:outside}).  The residual class---%
$z \notin P$ and $P \cap \Z = 0$---contains $(0)$ and all of
$\mathfrak{P}_1$, but not only these: for each rational prime $p$,
$p$-adic evaluation at $-p$ yields a surjection $\mRR \twoheadrightarrow
\Z_p$ whose kernel $Q_p$ is a prime lying outside $\mathfrak{P}_1$ and
distinct from every ideal above (Section~\ref{sec:outside}).  For primes
\emph{inside} $\mathfrak{P}_1$ we attach an analytic invariant
$\mathcal{Z}(P)$---the family of zero divisors of the unit-constant-term
elements of $P$---which satisfies a partition property; its completeness
is posed as a conjecture (Section~\ref{sec:invariant}).  We realize every
ultrafilter on a real-supported admissible divisor as the trace of a
prime ideal built by an ultraproduct construction
(Section~\ref{sec:ultra}).  Finally we treat the point-evaluation ideals
$P_a = \ker(\mathrm{ev}_a)$, proving primeness for all $a \in \D$ and
maximality for real nonzero $a$ (and for all nonzero $a$ in the Gaussian
ring $\mR$), and deduce a conditional B\'ezout theorem for functions with
disjoint divisors and coprime constant terms
(Sections~\ref{sec:pointeval}--\ref{sec:bezout}).

\medskip\noindent\textbf{Formalization.}
The trichotomy of maximal ideals, the classification of primes containing
$z$ or a nonzero integer, the unit-quotient lemma, part~(i) of
Proposition~\ref{prop:interesting} (in the strengthened form of
Remark~\ref{rem:onlymax}), the point-evaluation results, the
$p$-adic evaluation primes of Proposition~\ref{prop:padic} (including the
surjectivity of $\psi_p$ and the isomorphism $\mRR/Q_p \cong \Z_p$), the
partition property of Proposition~\ref{prop:prime1ultrafilter}, the
ultraproduct construction of Proposition~\ref{prop:ultraproductprime}
(in the real-supported form stated here), and the conditional B\'ezout
theorem (including the necessity of its coprimality hypothesis) have been
formally verified in Lean~4, building on the verified development of
\cite{PaperI}.  Every numbered theorem, proposition, lemma, and corollary
of this paper has been so verified; the conjecture and the questions of
Section~\ref{sec:open} have not.

Throughout, the arguments are stated for $\mRR$ for concreteness; they
apply verbatim to $\mR$ with $\Zi$ in place of $\Z$ unless noted.  We
write $\mathfrak{n}_0 = \{g : g(0) = 0\}$ for the augmentation ideal, and
$[z^n]f$ for the coefficient of $z^n$ in $f$.

\section{Notation and results recalled from \cite{PaperI}}
\label{sec:recall}

We recall the statements used below.  Proofs are in \cite{PaperI}; the
locators are placeholders to be synchronized with the final numbering of
the companion paper.

\begin{theorem}[Integer realization; {\cite[Thm.~1.1]{PaperI}}]
\label{thm:main}
An effective divisor $D$ on $\D$ is the zero divisor of a holomorphic
function on $\D$ with Taylor coefficients in $\Z$ if and only if $D$ is
invariant under complex conjugation.
\end{theorem}

\begin{theorem}[Gaussian-integer realization; {\cite[Thm.~1.2]{PaperI}}]
\label{thm:Zi}
Every effective divisor on $\D$ is the zero divisor of a holomorphic
function on $\D$ with Taylor coefficients in $\Zi$.
\end{theorem}

By an \emph{admissible divisor} on $\D$ we mean an effective divisor
(a multiplicity function on $\D$ with values in the nonnegative
integers) whose support meets every compact subset of $\D$ in a finite
set; the admissible divisors are exactly the zero divisors of functions
holomorphic on $\D$ and not identically zero.  In particular
(Theorem~\ref{thm:main}), every conjugation-invariant admissible
divisor in $\D$ assigning multiplicity $0$ to the origin is
the zero set, with multiplicity, of some element of $\mRR$ with constant
term $1$; this is the only input to the partition and ultraproduct
arguments below.  Both restrictions are necessary: the coefficients being
real forces the zero set to be conjugation-invariant, and constant term
$1$ forces nonvanishing at $0$.

\begin{proposition}[Units; {\cite[Prop.~4.1]{PaperI}}]
\label{prop:units}
An element $f \in \mR$ is a unit if and only if $f(0) \in \Zi^\times =
\{\pm 1, \pm i\}$ and $f$ is nowhere vanishing on $\D$.  The analogous
statement holds for $\mRR$, with $\Zi^\times$ replaced by $\Z^\times =
\{\pm 1\}$.
\end{proposition}

We also use, from \cite{PaperI}, that every $f \in \OD$ factors as
$f = g\,u$ with $g \in \mR$ and with $u$ a unit of $\OD$, that
consequently the
contraction $\phi \colon \mathrm{MaxSpec}(\OD) \to \mathrm{Spec}(\mR)$,
$\mathfrak{m} \mapsto \mathfrak{m} \cap \mR$, is injective, and that
$\mR/\mathfrak{n}_0 \cong \Zi$ (respectively $\mRR/\mathfrak{n}_0 \cong
\Z$).  These frame the picture but are not needed in the proofs that
follow.

\section{A trichotomy for maximal ideals}
\label{sec:trichotomy}

\begin{proposition}[Trichotomy of maximal ideals]
\label{prop:trichotomy}
Every maximal ideal $M \subset \mRR$ falls into exactly one of the
following types:
\begin{enumerate}
\item[(i)] $z \in M$,
\item[(ii)] $n \in M$ for some integer $n$ with $\abs{n} > 1$, but
  $z \notin M$,
\item[(iii)] neither $z$ nor any integer constant belongs to $M$.
\end{enumerate}
Moreover, type~\emph{(ii)} cannot occur: no maximal ideal of $\mRR$
contains an integer but not $z$.
\end{proposition}

\begin{proof}
The three types are mutually exclusive and exhaustive by definition.
It remains to show type~(ii) cannot occur.  Suppose $n \in M$ with
$\abs{n} > 1$.  Factoring $\abs{n}$ into rational primes inside $\mRR$
and using primeness of $M$, some prime factor $p$ lies in $M$.  Thus
$(p) \subseteq M$, and $M$ corresponds to a maximal ideal $\bar M$ of
$\mRR/(p) \cong \mathbb{F}_p[[z]]$ (the coefficientwise reduction is a
surjection with kernel $(p)$; surjectivity holds because any $\sum \bar
a_n z^n$ lifts to $\sum a_n z^n$ with $a_n \in \{0,\ldots,p-1\}$, of
radius of convergence $\geq 1$).  The ring $\mathbb{F}_p[[z]]$ is local
with maximal ideal $(\bar z)$, so $\bar M = (\bar z) \ni \bar z$, and
since $(p) \subseteq M$ this pulls back to $z \in M$.  Hence a maximal
ideal containing an integer of modulus exceeding $1$ contains $z$, and
type~(ii) is empty.
\end{proof}

\begin{definition}
A maximal ideal of type~(iii) is called a \emph{type~(iii) maximal
ideal}.
\end{definition}

We fix, once and for all, three classes of prime ideals of $\mRR$,
using fraktur notation:
\begin{itemize}
\item $\mathfrak{P} = \mathrm{Spec}(\mRR)$: all prime ideals of $\mRR$.
\item $\mathfrak{P}_1$: prime ideals $P$ containing at least one element
  of the form $1 + zF$ (i.e.\ with unit constant term).
\item $\mathfrak{M}$: maximal ideals of $\mRR$.
\end{itemize}
These satisfy $\mathfrak{M} \cap \mathfrak{P}_1 \subseteq \mathfrak{P}_1
\subseteq \mathfrak{P}$.  Every type~(iii) maximal ideal lies in
$\mathfrak{P}_1$ by Proposition~\ref{prop:interesting}(i) below;
type~(i) maximal ideals (containing $z$) do not contain
unit-constant-term elements, since if $z \in P$ and $1 + zF \in P$ then
$1 = (1 + zF) - z \cdot F \in P$, contradicting $P$ proper.  Whether
every element of $\mathfrak{P}_1$ is maximal is open
(Section~\ref{sec:open}).

\section{Prime ideals outside $\mathfrak{P}_1$}
\label{sec:outside}

\begin{proposition}[Primes containing $z$ or a nonzero integer]
\label{prop:primes_outside}
Every prime ideal $P$ of $\mRR$ falls into exactly one of the following
classes:
\begin{enumerate}
\item[(a)] $z \in P$, and then $P = (z)$ or $P = (z, p)$ for a rational
  prime $p$.  The ideal $(z)$ is prime and not maximal; the ideals
  $(z,p)$ are maximal, and are exactly the type~(i) maximal ideals.
\item[(b)] $z \notin P$ and $P \cap \Z \neq 0$, and then $P = (p)$ for a
  rational prime $p$.  These are prime but not maximal.
\item[(c)] $z \notin P$ and $P \cap \Z = 0$.
\end{enumerate}
Class~(c) contains $(0)$ and every member of $\mathfrak{P}_1$.  Indeed,
let the prime $P$ contain $1 + zF$.  If $z \in P$ then $1 = (1+zF) -
zF \in P$, impossible.  If a nonzero integer lies in $P$, then some
rational prime $p \in P$ (factor and use primeness), so $(p) \subseteq
P$ and the image of $P$ in $\mRR/(p) \cong \mathbb{F}_p[[z]]$ is a
proper ideal containing the image of $1 + zF$; but that image has
constant term $\bar 1 \neq 0$ and is therefore a unit of the local ring
$\mathbb{F}_p[[z]]$, a contradiction.  Class~(c) is strictly larger than
$\{(0)\} \cup \mathfrak{P}_1$: see Proposition~\ref{prop:padic}.
\end{proposition}

\begin{proof}
The three classes are mutually exclusive and exhaustive by definition.

\medskip\noindent\emph{Class (a): $z \in P$.}  Since $(z) \subseteq P$,
the image of $P$ in $\mRR/(z) \cong \Z$ (the isomorphism is evaluation at
$0$, whose kernel is exactly $(z)$: an element with constant term $0$ is
divisible by $z$ in $\mRR$) is a prime ideal of $\Z$, hence $(0)$ or
$(p)$ for a rational prime $p$.  Pulling back along the quotient, $P =
(z)$ or $P = (z,p)$.  Since $\mRR/(z,p) \cong \Z/(p) = \mathbb{F}_p$ is a
field, $(z,p)$ is maximal; since $\mRR/(z) \cong \Z$ is not a field,
$(z)$ is not.  A type~(i) maximal ideal contains $z$, hence is $(z,p)$
by the above; conversely each $(z,p)$ is maximal and contains $z$.

\medskip\noindent\emph{Class (b): $z \notin P$, $P \cap \Z \neq 0$.}
Fix a nonzero $n \in P \cap \Z$.  Factoring $\abs{n}$ into rational
primes inside $\mRR$ and using primeness of $P$ (the unit $\pm 1$ cannot
lie in $P$), some rational prime $p$ lies in $P$.  Now $(p) \subseteq
P$, so the image $\bar P$ of $P$ in $\mRR/(p) \cong \mathbb{F}_p[[z]]$
(the reduction of Proposition~\ref{prop:trichotomy}) is a legitimate
prime ideal---the kernel of the quotient is contained in $P$, which is
what the argument requires.  The prime ideals of the discrete valuation
ring $\mathbb{F}_p[[z]]$ are $(0)$ and $(\bar z)$; since $z \notin P$ we
have $\bar z \notin \bar P$, so $\bar P = (0)$, i.e.\ $P \subseteq (p)$.
With $(p) \subseteq P$ this gives $P = (p)$.  Finally
$\mathbb{F}_p[[z]]$ is not a field, so $(p)$ is not maximal.
\end{proof}

Class~(c) contains, besides $(0)$ and the members of $\mathfrak{P}_1$
studied below, the kernels of $p$-adic evaluations.

\begin{proposition}[$p$-adic evaluation primes]
\label{prop:padic}
Let $p$ be a rational prime and let $\Z_p$ denote the $p$-adic integers.
The assignment
\[
  \psi_p \colon \mRR \longrightarrow \Z_p, \qquad
  \psi_p\Big(\sum_{n \geq 0} a_n z^n\Big) = \sum_{n \geq 0} a_n (-p)^n
\]
(convergent $p$-adically, since $\abs{a_n(-p)^n}_p \leq p^{-n} \to 0$) is
a surjective ring homomorphism.  Its kernel $Q_p = \ker \psi_p$
satisfies:
\begin{enumerate}
\item[(i)] $Q_p$ is prime and not maximal, with $\mRR/Q_p \cong \Z_p$;
\item[(ii)] $z + p \in Q_p$, so $Q_p \neq (0)$;
\item[(iii)] $z \notin Q_p$ and $Q_p \cap \Z = 0$, so $Q_p$ lies in
  class~\emph{(c)} of Proposition~\ref{prop:primes_outside};
\item[(iv)] $Q_p \notin \mathfrak{P}_1$: no element of $Q_p$ has constant
  term $\pm 1$.  Indeed every $f \in Q_p$ has $p \mid f(0)$.
\end{enumerate}
In particular $\mathfrak{P} \setminus \mathfrak{P}_1$ strictly contains
$\{(0), (z)\} \cup \{(p)\} \cup \{(z,p)\}$.
\end{proposition}

\begin{proof}
The series defining $\psi_p$ converges in the complete nonarchimedean
ring $\Z_p$ because its terms tend to $0$.  Additivity is clear;
multiplicativity is the Cauchy-product identity, valid since both factor
series converge absolutely in the nonarchimedean sense.  Well-definedness
rests on the fact that an element of $\mRR$ determines its Taylor
coefficients.

\emph{Surjectivity.}  Given $w \in \Z_p$, define digits $\varepsilon_n
\in \{0, \ldots, p-1\}$ and remainders $w_n \in \Z_p$ by $w_0 = w$,
$\varepsilon_n \equiv w_n \pmod{p}$, and $w_{n+1} = (w_n -
\varepsilon_n)/(-p)$.  By construction $w = \sum_{n < N} \varepsilon_n
(-p)^n + (-p)^N w_N$ for every $N$, and $\abs{(-p)^N w_N}_p \leq p^{-N}
\to 0$, so $w = \sum_n \varepsilon_n (-p)^n$.  The digit sequence is
bounded, so $f = \sum_n \varepsilon_n z^n$ has radius of convergence at
least $1$ and integer coefficients, i.e.\ $f \in \mRR$, and $\psi_p(f) =
w$.

(i) $\mRR / Q_p \cong \Z_p$ by surjectivity; $\Z_p$ is an integral
domain, so $Q_p$ is prime, and $\Z_p$ is not a field, so $Q_p$ is not
maximal.  (ii) $\psi_p(z + p) = -p + p = 0$.  (iii) $\psi_p(z) = -p \neq
0$ and $\psi_p(n) = n \neq 0$ for nonzero $n \in \Z$.  (iv) For $f =
\sum a_n z^n$, $\psi_p(f) \equiv a_0 \pmod{p}$, since every term with
$n \geq 1$ lies in $p\Z_p$.  Hence $f \in Q_p$ forces $p \mid a_0$; in
particular $a_0 \neq \pm 1$.
\end{proof}

\begin{remark}
\label{rem:padic}
The primes $Q_p$ are $p$-adic analogues of the ideals $P_a$ of
Section~\ref{sec:pointeval}: where $P_a$ is the kernel of
archimedean evaluation at a point of $\D$, $Q_p$ is the kernel of
$p$-adic evaluation at $-p$, a point of $p\Z_p$.  More generally, any
$\alpha \in p\Z_p$ yields a ring homomorphism $\mRR \to \Z_p$, $f
\mapsto \sum a_n \alpha^n$, whose kernel is prime whenever it is proper;
for $\alpha = -p$ the base-$(-p)$ digit expansion yields surjectivity and
so identifies the quotient.  The primes $Q_p$ rule out any classification
of $\mathfrak{P} \setminus \mathfrak{P}_1$ by the ideals $(p)$ and
$(z,p)$ alone.  An argument for such a classification would pass to the
image of $P$ in $\mathbb{F}_p[[z]]$; that step is valid only when the
kernel $(p)$ of the reduction is contained in $P$, which is available in
the proof of Proposition~\ref{prop:primes_outside}(b) but not in
general---and $Q_p$, whose elements all have constant term divisible by
$p$ while $p \notin Q_p$, witnesses the failure.  A description of
class~\emph{(c)} in full remains open (Section~\ref{sec:open}).
\end{remark}

\section{Structure of primes in $\mathfrak{P}_1$}
\label{sec:invariant}

We first record a consequence of the unit characterization from
\cite{PaperI}.

\begin{lemma}[Unit quotient lemma]
\label{lem:unitquotient}
If $f_1, f_2 \in \mRR$ both have constant term $1$ and the same zero set
$Z(f_1) = Z(f_2)$ with multiplicities, then $f_1/f_2 \in \mRR^\times$.
\end{lemma}

\begin{proof}
Write $f_2 = \sum_{n \geq 0} a_n z^n$ with $a_0 = 1$ and $a_n \in \Z$.
The ratio $h = f_1/f_2$ extends holomorphically to $\D$ (the zeros of
$f_2$ are exactly those of $f_1$ with equal multiplicities, so $h$ has
no poles), is nowhere vanishing (identical zero sets), and $h(0) = 1$.
Writing $h = \sum_{n \geq 0} b_n z^n$ and comparing coefficients in $f_1
= f_2 \cdot h$ gives $b_0 = 1$ and
\[
  b_n = [z^n]f_1 - \sum_{k=1}^n a_k b_{n-k}, \quad n \geq 1.
\]
Since $[z^n]f_1 \in \Z$, $a_k \in \Z$, and $b_0 \in \Z$, induction gives
$b_n \in \Z$ for all $n$.  Because $h$ is nowhere vanishing on $\D$, $1/h
\in \OD$, so $h \in \mRR$.  Being nowhere vanishing with $h(0) = 1 \in
\Z^\times$, $h$ is a unit of $\mRR$ by Proposition~\ref{prop:units}.
\end{proof}

\begin{proposition}[Structure of type~(iii) maximal ideals]
\label{prop:interesting}
Let $M$ be a type~(iii) maximal ideal of $\mRR$.  Then:
\begin{enumerate}
\item[(i)] $M$ contains an element of the form $1 + zF$.  Indeed the set
  $I = \{g(0) : g \in M\} \subset \Z$ is an ideal, and $I = \Z$.
\item[(ii)] An element $f = 1 + zF \in \mRR$ is a unit of $\mRR$ if and
  only if it is nowhere vanishing on $\D$; this is the $\mRR$ case of
  Proposition~\ref{prop:units}, the constant term $1 \in \Z^\times$
  being automatic.  In particular, elements of $M$ of the form $1 + zF$
  are non-units, hence vanish somewhere in $\D$.
\item[(iii)] Two elements of $\mRR$ of the form $1 + zF$ with the same
  zero set differ by a unit of $\mRR$ (Lemma~\ref{lem:unitquotient}).
\end{enumerate}
\end{proposition}

\begin{proof}
(i) The set $I = \{g(0) : g \in M\}$ is an ideal of $\Z$ (closed under
addition since $M$ is an additive subgroup, and under integer
multiplication since $\Z \subset \mRR$ and $M$ is an ideal).  Suppose $I
\neq \Z$.

If $I = (0)$: every element of $M$ has zero constant term, i.e.\ $M
\subseteq \ker(\mathrm{ev}_0) = (z)$ (an element with constant term $0$
is divisible by $z$ in $\mRR$).  Since $(z)$ is a proper ideal and $M$ is
maximal, $M = (z) \ni z$, contradicting type~(iii).

Otherwise $I \subseteq (p)$ for some rational prime $p$ (take any prime
divisor of a generator of $I$).  Then $M \subseteq K_p := \{f \in \mRR :
p \mid f(0)\}$.  But $K_p = (z, p)$: the inclusion $\supseteq$ is clear,
and if $p \mid f(0)$ then $f = z \cdot \tfrac{f - f(0)}{z} + p \cdot
\tfrac{f(0)}{p}$ exhibits $f \in (z,p)$, the first quotient lying in
$\mRR$ by divisibility of constant-term-zero elements by $z$.  Since
$K_p$ is proper ($1 \notin K_p$) and $M$ is maximal, $M = K_p = (z,p)
\ni z$, again contradicting type~(iii).

Hence $I = \Z$, and some element of $M$ has constant term $1$.

(ii) This is Proposition~\ref{prop:units} for $\mRR$ applied to $f = 1 +
zF$ (whose constant term is the unit $1$): $f$ is a unit iff it is
nowhere vanishing.  Elements of $M$ are non-units, so those of the form
$1 + zF$ must vanish somewhere in $\D$.

(iii) Immediate from Lemma~\ref{lem:unitquotient}.
\end{proof}

\begin{remark}
\label{rem:onlymax}
The proof of part~(i) uses only that $M$ is maximal and $z \notin M$; the
type~(iii) hypothesis that $M$ contains no nonzero integer is not needed
(and indeed follows from these two by
Proposition~\ref{prop:trichotomy}).  So every maximal ideal not
containing $z$ contains an element of constant term $1$; equivalently,
$\mathfrak{M} \cap \mathfrak{P}_1$ consists exactly of the maximal ideals
avoiding $z$.
\end{remark}

\begin{remark}
\label{rem:nogeneration}
Part~(i) does not imply that the elements of the form $1 + zF$ lying in
$M$ generate $M$ as an ideal.  It supplies \emph{some} element of $M$
with constant term $1$, not that every $z$-reduced element of $M$ has
unit constant term: removing the leading power of $z$ from $f \in M$
leaves an element of $M$ whose constant term can be any value in $I$.
Whether the unit-constant-term elements of $M$ generate $M$ is open, and
the question propagates: the completeness of the invariant $\mathcal{Z}$
below is stated as a conjecture because the natural proof requires this
generation property.
\end{remark}

\begin{definition}
For a type~(iii) maximal ideal $M$, define its \emph{analytic invariant}
\[
  \mathcal{Z}(M) = \{ Z(f) : f \in M,\; f = 1 + zF \text{ for some }
  F \in \mRR \},
\]
where $Z(f) \subset \D$ is the zero set of $f$ counted with
multiplicities.
\end{definition}

\begin{remark}
\label{rem:Zunion}
The family $\mathcal{Z}(M)$ is closed under finite unions: if $f, g \in
M$ then $fg \in M$ and $Z(fg) = Z(f) \cup Z(g)$.  Each member of
$\mathcal{Z}(M)$ is an admissible divisor which is conjugation-invariant
(the coefficients of its realizing function are real) and assigns
multiplicity $0$ to the origin (the constant term is $1$); by
Theorem~\ref{thm:main}, every divisor with these properties is $Z(h)$
for some $h \in \mRR$ with constant term $1$.
\end{remark}

The invariant extends to all of $\mathfrak{P}_1$.  For $P \in
\mathfrak{P}_1$, define $\mathcal{Z}(P)$ by the same formula.

\begin{conjecture}[Completeness of $\mathcal{Z}$]
\label{prop:prime1invariant}
For $P_1, P_2 \in \mathfrak{P}_1$: $P_1 = P_2 \iff \mathcal{Z}(P_1) =
\mathcal{Z}(P_2)$.  In particular type~(iii) maximal ideals would be
determined by their invariants.
\end{conjecture}

\begin{remark}
If the unit-constant-term elements of $P$ generate $P$ for every $P \in
\mathfrak{P}_1$ (the open generation property of
Remark~\ref{rem:nogeneration}), the conjecture follows: given
$\mathcal{Z}(P_1) = \mathcal{Z}(P_2)$, each unit-constant-term $f \in
P_1$ admits a unit-constant-term $g \in P_2$ with $Z(g) = Z(f)$, whence
$f = ug$ for a unit $u$ by Proposition~\ref{prop:interesting}(iii) and
$f \in P_2$; generation then gives $P_1 \subseteq P_2$, and symmetry
concludes.  The generation property itself resists the natural attack:
stripping powers of $z$ and reducing constant terms against a fixed $1 +
zG \in P$ produces, from $f \in P$, expressions $f = c \cdot g + z
\cdot h$ with $h \in P$, and iterating yields an infinite descent rather
than the finite $\mRR$-combination that ideal generation requires.
\end{remark}

\begin{proposition}[Partition property of $\mathcal{F}_Z(P)$]
\label{prop:prime1ultrafilter}
Let $P \in \mathfrak{P}_1$ and $f = 1 + zF \in P$ with $Z = Z(f)$.  Put
\[
  \mathcal{F}_Z(P) = \{ Z(h) : h \in P,\; h = 1 + zH,\; Z(h) \subseteq Z \}.
\]
For every partition $Z = Z_1 \sqcup Z_2$ into conjugation-invariant
subdivisors, at least one of $Z_1, Z_2$ lies in $\mathcal{F}_Z(P)$.
\end{proposition}

\begin{proof}
Note first that $Z = Z(f)$ is conjugation-invariant (the coefficients of
$f$ are real) and assigns multiplicity $0$ to the origin ($f(0) = 1$),
and these properties pass to the conjugation-invariant pieces $Z_1,
Z_2$.  By Theorem~\ref{thm:main}, obtain $f_1, f_2 \in \mRR$ with
constant term $1$ and $Z(f_i) = Z_i$.  Then $f_1 f_2$ has constant term
$1$ and $Z(f_1 f_2) = Z_1 \cup Z_2 = Z = Z(f)$, so by
Proposition~\ref{prop:interesting}(iii), $f_1 f_2 = u f$ for a unit $u$.
As $f \in P$, $f_1 f_2 \in P$; primeness gives $f_1 \in P$ or $f_2 \in
P$, i.e.\ $Z_1 \in \mathcal{F}_Z(P)$ or $Z_2 \in \mathcal{F}_Z(P)$.
\end{proof}

\begin{remark}
The restriction to conjugation-invariant pieces is necessary for the
proof, since only conjugation-invariant divisors are realized over $\Z$
(Theorem~\ref{thm:main}).  When $|Z|
\subseteq \R$ every subdivisor is conjugation-invariant and the property
holds for arbitrary partitions.
\end{remark}

\begin{remark}
The partition property is all that follows from primeness alone.  The
full ultrafilter property---that \emph{exactly} one of $Z_1, Z_2$ lies in
$\mathcal{F}_Z(P)$---requires that $P$ not contain elements with disjoint
zero sets, which depends on maximality or further structural information
about $P$.
\end{remark}

\section{Ultrafilters give prime ideals via ultraproducts}
\label{sec:ultra}

\begin{proposition}[Ultraproduct construction]
\label{prop:ultraproductprime}
Let $Z$ be an admissible divisor in $\D$ with $|Z| \subseteq \R
\setminus \{0\}$, and let $\mathcal{U}$ be an ultrafilter on the
underlying point set $|Z|$ (ignoring multiplicities).  For each $W
\subseteq |Z|$, the set $W$ is a conjugation-invariant admissible
divisor avoiding the origin, so Theorem~\ref{thm:main} supplies $f_W \in
\mRR$ with constant term $1$ and $Z(f_W) = W$ as a set.  Define
\[
  \Phi_\mathcal{U} \colon \mRR \to \prod_{a \in Z} \C \Big/ \mathcal{U},
  \qquad \Phi_\mathcal{U}(h) = [(h(a))_{a \in Z}]_\mathcal{U}.
\]
Then:
\begin{enumerate}
\item[(i)] $\Phi_\mathcal{U}$ is a ring homomorphism to a field.  Here
  $\prod_{a \in Z} \C / \mathcal{U}$ is the ultraproduct of copies of
  $\C$ along $\mathcal{U}$; since $\mathcal{U}$ is a proper ultrafilter
  and each factor is a field, {\L}o\'s's theorem \cite{CK} implies the
  ultraproduct is a field.
\item[(ii)] $P_\mathcal{U} = \ker(\Phi_\mathcal{U}) \in \mathfrak{P}_1$:
  it is prime (kernel of a map to a field) and contains $f_Z$, which has
  constant term $1$.
\item[(iii)] $\mathcal{F}_Z(P_\mathcal{U}) = \mathcal{U}$: for $h = 1 +
  zH$ with $Z(h) \subseteq Z$, one has $h \in P_\mathcal{U}$ iff $\{a
  \in Z : h(a) = 0\} = Z(h) \in \mathcal{U}$.
\item[(iv)] The map $\mathcal{U} \mapsto P_\mathcal{U}$ is injective: if
  $P_{\mathcal{U}_1} = P_{\mathcal{U}_2}$ then $\mathcal{U}_1 =
  \mathcal{F}_Z(P_{\mathcal{U}_1}) = \mathcal{F}_Z(P_{\mathcal{U}_2}) =
  \mathcal{U}_2$.
\end{enumerate}
\end{proposition}

\begin{proof}
(i) The reduced product $\prod_{a \in Z}\C/\mathcal{U}$ is the
ultraproduct of copies of $\C$ along the proper ultrafilter
$\mathcal{U}$; by {\L}o\'s's theorem for the first-order theory of fields
it is a field, and $\Phi_\mathcal{U}$ is a ring homomorphism by
coordinatewise evaluation.  (No model theory is actually required here:
a nonzero class is represented by a family nonvanishing on a member of
$\mathcal{U}$, and inverting it there and setting it to $1$ elsewhere
produces an inverse class.)  (ii) The kernel of a homomorphism to a
field is prime; $f_Z \in P_\mathcal{U}$ since $Z \in \mathcal{U}$ (the
whole index set belongs to every filter) and $f_Z$ vanishes on $Z$.
(iii) For $h = 1 + zH$ with $Z(h) \subseteq Z$: by definition of the
reduced product, $h \in P_\mathcal{U}$ iff $\{a \in Z : h(a) = 0\} \in
\mathcal{U}$, and the restriction $Z(h) \subseteq Z$ makes that set
equal $Z(h)$; conversely every $W \in \mathcal{U}$ arises as $Z(f_W)$
with $f_W \in P_\mathcal{U}$ by the membership criterion just
established.  (iv) Immediate from (iii).
\end{proof}

\begin{remark}
The map $\mathcal{U} \mapsto P_\mathcal{U}$ from ultrafilters on
real-supported admissible divisors into $\mathfrak{P}_1$ is injective by
part~(iv).  Whether it is surjective---whether every element of
$\mathfrak{P}_1$ arises as some $P_\mathcal{U}$---is open
(Section~\ref{sec:open}).  The restriction $|Z| \subseteq \R$ enters
only through the realizability of arbitrary subsets $W \subseteq |Z|$
over $\Z$; for general $Z$ the same construction works verbatim provided
$\mathcal{U}$ admits a base of conjugation-invariant sets, and works
over the Gaussian ring $\mR$ with no restriction beyond $0 \notin |Z|$
(Theorem~\ref{thm:Zi}).
\end{remark}

\section{Point-evaluation ideals}
\label{sec:pointeval}

For $a \in \D$ let $\mathrm{ev}_a \colon \mRR \to \C$, $f \mapsto f(a)$,
and $P_a = \ker(\mathrm{ev}_a)$.

\begin{proposition}[Primeness]
\label{prop:pointeval_prime}
For each $a \in \D$, the ideal $P_a$ is prime in $\mRR$.
\end{proposition}

\begin{proof}
Since every $f \in \mRR$ is holomorphic on $\D$ and $a \in \D$,
$\mathrm{ev}_a$ is a ring homomorphism, nonzero (the constant $1 \mapsto
1$) with codomain the integral domain $\C$; hence its kernel $P_a$ is
prime.
\end{proof}

\begin{proposition}[Maximality for real points]
\label{prop:pointeval_max_real}
For each $a \in \R \cap \D$, $P_a$ is a maximal ideal of $\mRR$ with
$\mRR/P_a \cong \R$.
\end{proposition}

\begin{proof}
Since all coefficients are real integers, $\mathrm{ev}_a(\mRR) \subseteq
\R$ for $a \in \R$.  The case $a = 0$ is excluded ($P_0 =
\mathfrak{n}_0$, with $\mRR/P_0 \cong \Z$ not a field; indeed $z \in
P_0$, so $P_0$ is type~(i)).  For $a \in (-1,1) \setminus \{0\}$ and any
$w \in \R$, define remainders $r_n = w - \sum_{k=0}^n c_k a^k$ by $r_{-1}
= w$ and, at step $n$, $c_n \in \Z$ the nearest integer to
$r_{n-1}/a^n$, so $r_n = r_{n-1} - c_n a^n$ and $\abs{r_n} \leq
\tfrac12\abs{a}^n$.

The bound $\abs{a} < 1$ keeps $\{c_n\}$ bounded: for $n \geq 1$,
$\abs{r_{n-1}} \leq \tfrac12\abs{a}^{n-1}$, so $\abs{r_{n-1}/a^n} \leq
\tfrac{1}{2\abs{a}}$ and $\abs{c_n} \leq \tfrac{1}{2\abs{a}} + \tfrac12$,
a constant independent of $n$.  Hence $\sup_n \abs{c_n} < \infty$, so
$\limsup \abs{c_n}^{1/n} = 1$ and $f = \sum_{n \geq 0} c_n z^n \in \mRR$.
Since $\abs{r_N} \leq \tfrac12\abs{a}^N \to 0$, the partial sums
$\sum_{k=0}^N c_k a^k = w - r_N \to w$, so $f(a) = w$.  Thus
$\mathrm{ev}_a \colon \mRR \twoheadrightarrow \R$ is surjective, $\mRR/P_a
\cong \R$ is a field, and $P_a$ is maximal.
\end{proof}

\begin{proposition}[Maximality in $\mR$]
\label{prop:pointeval_max_gaussian}
For each nonzero $a \in \D$, $P_a = \ker(\mathrm{ev}_a)$ is a maximal
ideal of $\mR = \Zi[[z]] \cap \OD$, with $\mR/P_a \cong \C$.
\end{proposition}

\begin{proof}
The case $a = 0$ is excluded: $P_0 = \mathfrak{n}_0$ with
$\mR/P_0 \cong \Zi$, not a field.  For $a \neq 0$, the same bounded radix
expansion, with $c_n \in \Zi$ chosen nearest to $r_{n-1}/a^n$, gives
$\abs{r_n} \leq \tfrac{\sqrt2}{2}\abs{a}^n$ and $\abs{c_n} \leq
\tfrac{\sqrt2}{2\abs{a}} + \tfrac{\sqrt2}{2}$, independent of $n$; so $f
= \sum c_n z^n \in \mR$ with $f(a) = w$ for any $w \in \C$.  Hence
$\mR/P_a \cong \C$ and $P_a$ is maximal.
\end{proof}

\begin{remark}
Whether $P_a$ is maximal in $\mRR$ for $a \in \D \setminus \R$ remains
open: $\mathrm{ev}_a(\mRR)$ is a subring of $\C$ containing $\Z[a]$, but
surjectivity onto $\C$ with only integer (not Gaussian-integer)
coefficients is an arithmetic question about $a$.
\end{remark}

\section{A conditional B\'ezout property}
\label{sec:bezout}

\begin{theorem}[Conditional B\'ezout property]
\label{thm:bezout}
Let $f, g \in \mRR$ have disjoint zero sets in $\D$ and coprime constant
terms, $\gcd(f(0), g(0)) = 1$.  Then $(f,g)$ is not contained in any
type~(i) maximal ideal, nor in any point-evaluation ideal $P_a$ with $a
\in \D$.  Consequently, if every type~(iii) maximal ideal is a
point-evaluation ideal $P_a$ for some $a \in \D$, then $(f,g) = \mRR$.
\end{theorem}

\begin{proof}
Suppose $(f,g) \subseteq M$ for a maximal ideal $M$.  By
Proposition~\ref{prop:trichotomy}, $M$ is type~(i) or type~(iii).

\emph{Type~(i):} by Proposition~\ref{prop:primes_outside}(a), $M =
(z,p)$ for a rational prime $p$.  An element $zA + pB \in (z,p)$ has
constant term $p \cdot B(0)$, so membership of $f$ and $g$ in $M$ forces
$p \mid f(0)$ and $p \mid g(0)$, contradicting $\gcd(f(0), g(0)) = 1$.

\emph{Type~(iii), $M = P_a$ with $a \in \D$:} then $f(a) = g(a) = 0$, so
$a \in Z(f) \cap Z(g)$, contradicting disjointness of the zero sets.

If every type~(iii) maximal ideal is such a $P_a$, then $(f,g)$ lies in
no maximal ideal, hence $(f,g) = \mRR$.
\end{proof}

\begin{remark}[Sharpness]
\label{rem:bezoutsharp}
Coprimality of the constant terms cannot be dropped.  For any
rational prime $p$, the pair $f = g = p$ has empty (hence disjoint) zero
sets in $\D$,
yet $(f, g) = (p) \subseteq (z,p)$ is proper.  The same pair shows that
membership of $f$ in a type~(i) maximal ideal does not force $f(0) = 0$:
only divisibility of $f(0)$ by $p$ follows.
\end{remark}

\begin{remark}
The conclusion is non-constructive even conditionally: Zorn's lemma
produces the maximal ideal and no explicit cofactors are exhibited.
Whether every type~(iii) maximal ideal is a point-evaluation ideal is the
key open question; see Section~\ref{sec:open}.
\end{remark}

\begin{remark}[Partial structure of $\mathrm{Spec}(\mRR)$]
\label{rem:spec}
Assembling the above: the primes containing $z$ are $(z)$ and the
maximal $(z,p)$; the primes containing a nonzero integer but not $z$ are
the non-maximal $(p)$ (Proposition~\ref{prop:primes_outside}).  The
remaining primes contain neither, and include $(0)$; the $p$-adic
evaluation kernels $Q_p \cong \ker(\mRR \twoheadrightarrow \Z_p)$, which
lie outside $\mathfrak{P}_1$ (Proposition~\ref{prop:padic}); and, inside
$\mathfrak{P}_1$, the point-evaluation ideals $P_a$---maximal for
nonzero $a \in \R \cap \D$
(Proposition~\ref{prop:pointeval_max_real}), prime for all $a \in \D$
(Proposition~\ref{prop:pointeval_prime})---and the kernels
$P_\mathcal{U}$ of ultraproduct maps, with $\mathcal{F}_Z(P_\mathcal{U})
= \mathcal{U}$ (Proposition~\ref{prop:ultraproductprime}).  Whether every
type~(iii) maximal ideal is a $P_a$, whether $P_a$ is maximal for
non-real $a$, and whether every element of $\mathfrak{P}_1$ is maximal
are open.
\end{remark}

\section{Further Questions}
\label{sec:open}

\begin{itemize}
\item \textbf{Classification of type~(iii) maximal ideals.}  The
  point-evaluation ideals $P_a$ ($a \in \D$) are type~(iii) maximal
  ideals.  The central open question is whether every type~(iii)
  maximal ideal has this form, i.e.\ whether $\mathfrak{M} \cap
  \mathfrak{P}_1$ consists exactly of the $P_a$.  A
  positive answer yields the full B\'ezout property
  (Theorem~\ref{thm:bezout}) unconditionally.

\item \textbf{The generation property.}  Is every $P \in \mathfrak{P}_1$
  generated, as an ideal, by its elements of the form $1 + zF$
  (Remark~\ref{rem:nogeneration})?  A positive answer proves
  Conjecture~\ref{prop:prime1invariant}, the completeness of the
  invariant $\mathcal{Z}$.

\item \textbf{Structure of class (c).}  Classify the primes containing
  neither $z$ nor a nonzero integer
  (Proposition~\ref{prop:primes_outside}(c)).  Besides $(0)$ and
  $\mathfrak{P}_1$, the class contains the $p$-adic evaluation primes
  $Q_p$ (Proposition~\ref{prop:padic}) and, more generally, kernels of
  evaluations $f \mapsto \sum a_n \alpha^n$ at $\alpha \in p\Z_p$.  Do
  archimedean and $p$-adic evaluations (and their ultraproducts)
  account for all of class~(c)?  Is $Q_p$ principal---indeed, is $Q_p =
  (z + p)$?

\item \textbf{Surjectivity of the ultraproduct map.}  The injection
  $\mathcal{U} \mapsto P_\mathcal{U}$
  (Proposition~\ref{prop:ultraproductprime}) parametrizes a family of
  primes in $\mathfrak{P}_1$ by ultrafilters on real-supported
  admissible divisors.  Whether it is surjective onto $\mathfrak{P}_1$,
  and whether the $P_\mathcal{U}$ are maximal, is open.

\item \textbf{Maximality inside $\mathfrak{P}_1$.}  Is every element of
  $\mathfrak{P}_1$ maximal?  Is $P_a$ maximal in $\mRR$ for non-real $a
  \in \D$?

\item \textbf{Surjectivity of $\phi$.}  Whether every maximal ideal of
  $\mR$ not containing $\mathfrak{n}_0$ arises by contraction from $\OD$
  is open; this would require $\mathfrak{n} \cdot \OD$ proper for every
  such maximal $\mathfrak{n} \subset \mR$.  (The contraction $\phi$ and
  its injectivity are from \cite{PaperI}.)

\item \textbf{Finer unit structure.}  Proposition~\ref{prop:units}
  characterizes units; the finer structure of $\mR^\times$ and
  $\mRR^\times$ is discussed in \cite{PaperI} and not revisited here.
\end{itemize}

\begin{thebibliography}{9}

\bibitem{PaperI}
J.~Bannon and D.~V.~Feldman,
\textit{Integer Coefficients Power Series with Prescribed Zero Sets},
companion paper.

\bibitem{CK}
C.~C.~Chang and H.~J.~Keisler,
\textit{Model Theory}, 3rd ed.,
North-Holland, Amsterdam, 1990.

\end{thebibliography}

\end{document}
